% Part: first-order-logic
% Chapter: sequent-calculus
% Section: provability-quantifiers

\documentclass[../../../include/open-logic-section]{subfiles}

\begin{document}

\olfileid{fol}{seq}{qpr}
\olsection{可推导性与量词}

\begin{explain}
  完备性定理也要求相继式演算规则能够得出本节所建立的关于~$\Proves$ 的事实。
\end{explain}

\begin{thm}
\ollabel{thm:strong-generalization} 如果 $c$ 是不在 $\Gamma$ 或 $!A(x)$ 中出现的常项，且 $\Gamma \Proves !A(c)$，则 $\Gamma
\Proves \lforall[x][!A(x)]$。
\end{thm}

\begin{proof}
令 $\pi_0$ 为 $\Gamma_0 \Sequent !A(c)$ 的一个 $\Log{LK}$-推导，其中 $\Gamma_0$ 是 $\Gamma$ 的某个有限子集。通过增加一个 $\RightR{\lforall}$ 推理，我们得到 $\Gamma_0 \Sequent
\lforall[x][!A(x)]$ 的一个推导，因为 $c$ 不在 $\Gamma$ 或 $!A(x)$ 中出现，因而满足特征变元条件。
\end{proof}

\begin{prop}
\ollabel{prop:provability-quantifiers}
\begin{tagenumerate}{prvEx,prvAll}
\tagitem{prvEx}{$!A(t) \Proves \lexists[x][!A(x)]$.}{}
\tagitem{prvAll}{$\lforall[x][!A(x)] \Proves !A(t)$.}{}
\end{tagenumerate}
\end{prop}

\begin{proof}
\begin{tagenumerate}{prvEx,prvAll}
\tagitem{prvEx}{相继式 $!A(t) \Sequent \lexists[x][!A(x)]$ 是可推导的：
  \begin{prooftree}
    \Axiom$!A(t) \fCenter !A(t)$
    \RightLabel{\RightR{\lexists}}
    \UnaryInf$!A(t) \fCenter \lexists[x][!A(x)]$
    \end{prooftree}}{}
\tagitem{prvAll}{相继式 $\lforall[x][!A(x)] \Sequent !A(t)$ 是可推导的：
  \begin{prooftree}
    \Axiom$!A(t) \fCenter !A(t)$
    \RightLabel{\LeftR{\lforall}}
    \UnaryInf$\lforall[x][!A(x)] \fCenter !A(t)$
    \end{prooftree}}{}
\end{tagenumerate}
\end{proof}

\end{document}